001|heading 1|||Initial Project Plan for Schedify: Constraint-Based Automated Exam Scheduling System for University Departments
002||||Date: [07/03/2025]
003||||Prepared by: 
004|heading 3|||Team #18
005||||Emre Demirbaş - 21050111069
006||||Ömer Faruk Özüyağlı - 21050111047
007||||Mehmet Ali Yılmaz - 21050111057
008|heading 3|||1. Introduction
009|heading 4|||Project Overview
010||||The Constraint-Based Automated Exam Scheduling System, named Schedify, is a software solution designed to automate the complex and time-consuming task of exam scheduling for universities. The system will ensure optimal utilization of resources, such as classrooms and instructors, while respecting to various constraints, including classroom capacities, instructor availability, and student exam conflicts. By using advanced scheduling algorithms, Schedify aims to generate conflict-free, efficient, and fair exam schedules for multiple departments within a university.
011|heading 4|||Business Goals
012||||The primary objective of Schedify is to streamline the exam scheduling process, reducing manual effort and minimizing errors caused by humans. The system will provide the following benefits:
013||6|0|Efficiency: Automate the scheduling process to save time and resources for university administration.
014||6|0|Fairness: Ensure conflict-free schedules for students and instructors while respecting constraints.
015||6|0|Scalability: Support universities of varying sizes, from small institutions to large multi-depth The Constraint-Based Automated Exam Scheduling System, named Schedify, is a software solution designed to automate the complex and time-consuming task of exam scheduling for universities. The system will ensure optimal utilization of resources, such as classrooms and instructors, while adhering to various constraints, including classroom capacities, instructor availability, and student exam conflicts. By leveraging advanced scheduling algorithms, Schedify aims to generate conflict-free, efficient, and fair exam schedules for multiple departments within a university.
016||6|0|Improved Communication: Provide clear and timely notifications to stakeholders, including students, faculty, and administrators, about finalized schedules and changes.
017||||Schedify will ultimately enhance the operational efficiency of universities, improve stakeholder satisfaction, and ensure compliance with institutional policies and constraints.
019|heading 3|||2. Project Organization
020|heading 4|||Team Members & Roles
021||||The development of the Schedify: Constraint-Based Automated Exam Scheduling System will involve a multidisciplinary team with clearly defined roles to ensure smooth execution and collaboration. The team structure is as follows:
022||4|0|Project Manager: Ömer Faruk Özüyağlı Responsible for overseeing the project, managing timelines, coordinating team efforts, and ensuring the project meets its objectives.
023||4|0|Lead Developer: Emre Demirbaş Oversees the software architecture, ensures adherence to coding standards, and leads the development team.
024||4|0|Backend Developers: Mehmet Ali Yılmaz, Emre Demirbaş Focus on implementing the scheduling engine, database management, and APIs for system functionality.
025||4|0|Frontend Developers: Emre Demirbaş, Ömer Faruk Özüyağlı Responsible for designing and developing the user interface, ensuring usability and accessibility.
026||4|0|UI/UX Designers: Mehmet Ali Yılmaz Create intuitive and user-friendly designs, including mockups and prototypes for the system.
027||4|0|Quality Assurance (QA) Team: Ömer Faruk ÖzüyağlıConduct testing to identify and resolve bugs, ensure system reliability, and validate that the system meets requirements.
028||4|0|System Administrator: Mehmet Ali Yılmaz Manages deployment, server configurations, and system maintenance.
029|heading 4|||Stakeholders
030||||The key stakeholders for the project include:
031||7|0|University Administration:
032||7|1|Registrar and Exam Coordinators: Oversee the scheduling process and provide data inputs.
033||7|1|IT Department: Provide technical support and integration with existing university systems.
034||7|0|Academic Departments:
035||7|1|Department Heads: Define department-specific constraints and requirements.
036||7|1|Faculty Members: Provide availability and preferences for exam schedules.
037||7|0|Students and Instructors: End-users of the system who will access their exam schedules.
038||7|0|Development Team: The internal team responsible for designing, developing, and deploying the system.
039||7|0|External Advisors: Professors or industry experts providing guidance on best practices and technical challenges.
040||||This organizational structure ensures that all aspects of the project are covered, from development to deployment, while maintaining clear communication with stakeholders.
048|heading 3|||3. Scope & Deliverables
049|heading 4|||Scope Definition
050||||The Schedify: Constraint-Based Automated Exam Scheduling System will focus on automating the exam scheduling process for universities while respecting to various constraints. The scope of the project includes the following:
051||||In-Scope:
052||5|0|Development of a user management system with role-based access control for administrators, department coordinators, faculty, and students.
053||5|0|Implementation of a data management module to handle courses, students, instructors, and classrooms.
054||5|0|Integration of a constraint management system to enforce rules such as classroom capacity, instructor availability, and student exam conflict avoidance.
055||5|0|Development of an automated exam scheduling algorithm capable of generating conflict-free schedules based on predefined constraints.
056||5|0|Support for manual adjustments to schedules by authorized users.
057||5|0|Creation of a reporting and notification system to generate detailed reports and notify stakeholders of finalized schedules or changes.
058||5|0|Development of a user-friendly interface for schedule visualization and management.
059||5|0|Integration with existing university systems for data import/export.
061|heading 4|||Key Deliverables
062||||The project will produce the following deliverables:
063||10|0|Software Requirements Document (SRD): A detailed document outlining functional and non-functional requirements, constraints, and assumptions.
064||10|0|System Design Artifacts:
065||10|1|Architectural diagrams, including system architecture and database schema.
066||10|1|UI/UX mockups and prototypes for the user interface.
067||10|0|Prototype UI: A working prototype of the user interface showcasing core functionalities such as schedule visualization and basic user interactions.
068||10|0|Automated Scheduling Engine: A fully functional scheduling engine capable of generating conflict-free exam schedules based on constraints.
069||10|0|Fully Functional Software: The complete system, including all modules (user management, data management, constraint management, scheduling engine, reporting, and notifications).
070||10|0|Testing Artifacts:
071||10|1|Test cases and test reports for functional, performance, and user acceptance testing.
072||10|1|Bug tracking and resolution documentation.
073||10|0|User Documentation:
074||10|1|User manuals for administrators, department coordinators, and faculty.
075||10|1|Training materials for end-users.
076||10|0|Deployment Package:
077||10|1|Deployed system on a cloud platform or university server.
078||10|1|Technical documentation for system administrators, including deployment and maintenance instructions.
079||||By clearly defining the scope and deliverables, the project ensures alignment with stakeholder expectations and provides a roadmap for successful implementation.
081|heading 2|||4. Work Breakdown Structure (WBS)
103|heading 2|||5. Project Schedule
104||||• Start Date: [25/03/2025]
105||||• Estimated Completion Date: [10/06/2025]
106||||• Milestones:
107||||  - M1: Requirements Approval
108||||  - M2:System Design Approval
109||||  - M3: Core Backend Completed
110||||  - M4: Frontend Development Complete
111||||  - M5: System Integration & Testing Complete
112||||  - M6: Final Deployment & Documentation
114|heading 2|||6. Risk Analysis & Mitigation
118|heading 2|||7. Resource Requirements
120|heading 4|||Hardware Resources
121||2|0|Cloud Infrastructure:
122||2|1|Cloud servers for hosting the application (e.g., AWS, Google Cloud, or Azure).
123||2|1|Scalable virtual machines to handle peak loads during scheduling periods.
124||2|0|Testing Devices:
125||2|1|Laptops and desktops for development and testing.
126||2|1|Mobile devices (optional) for testing responsiveness and compatibility.
128||||Software Resources
129||9|0|Development Tools:
130||9|1|Backend: Java (Spring Boot)
131||9|1|Frontend: React.js 
132||9|1|Database: PostgreSQL
133||9|0|Scheduling Algorithms:
134||9|1|Constraint Satisfaction Problem (CSP) solvers
135||9|0|Collaboration Tools:
136||9|1|Version control: Git and GitHub.
137||9|1|Task management: Jira for tracking progress.
138||9|0|Testing Tools:
139||9|1|JUnit automated testing framework.
140||9|1| JMeter load testing tool.
141||9|0|Deployment Tools:
142||9|1|Docker and Kubernetes for containerization and orchestration.
143||9|1|CI/CD pipelines for automated builds and deployments.
145|heading 4|||Personnel Resources
146||11|0|Project Manager:
147||11|1|Oversees the project timeline, resource allocation, and stakeholder communication.
148||11|0|Backend Developers:
149||11|1|Responsible for implementing the scheduling engine, APIs, and database integration.
150||11|0|Frontend Developers:
151||11|1|Develop the user interface and ensure usability and responsiveness.
152||11|0|UI/UX Designers:
153||11|1|Create intuitive designs and prototypes for the system.
154||11|0|Quality Assurance (QA) Team:
155||11|1|Conduct functional, performance, and user acceptance testing.
156||11|0|System Administrator:
157||11|1|Manages deployment, server configurations, and system maintenance.
159|heading 2|||8. Monitoring & Reporting
160|heading 2|||Monitoring Tools
162||1|0|Project Management Tools:
163||1|1|Jira: For task tracking, gantt chart creation, sprint planning, and backlog management. 
164||1|0|Version Control:
165||1|1|Git with GitHub: For source code management, version tracking, and collaboration.
166||1|0|Communication Tools:
167||1|1|Microsoft Teams: For team communication, updates for weekly progress meetings and stakeholder reviews.
168|heading 4|||Reporting Processes
169||3|0|Weekly Progress Meetings:
170||3|1|Conducted every Monday to review completed tasks, discuss blockers, and plan for the upcoming week.
171||3|1|Attendees: Project Manager, Developers, QA Team, and other relevant stakeholders.
172||3|0|Bi-Weekly Sprint Reviews:
173||3|1|At the end of each sprint, the team will present completed features and gather feedback from stakeholders.
174||3|1|Focus on ensuring deliverables meet requirements and identifying areas for improvement.
175||3|0|Monthly Status Reports:
176||3|1|A detailed report summarizing progress, risks, and budget utilization will be shared with stakeholders.
177||3|1|Includes updates on milestones, task completion rates, and any deviations from the schedule.
180|heading 4|||Key Performance Indicators (KPIs)
181||||To measure the project's progress and success, the following KPIs will be tracked:
182||13|0|Task Completion Rate: Percentage of completed tasks versus planned tasks for each sprint.
183||13|0|Milestone Achievement: Timely completion of key milestones (e.g., Requirements Approval, Backend Completion).
184||13|0|Bug Count: Number of bugs identified and resolved during testing phases.
185||13|0|Schedule Adherence: Percentage of tasks completed on time versus delayed tasks.
186||13|0|Stakeholder Feedback: Satisfaction levels from stakeholders during sprint reviews and user acceptance testing.
187|heading 4|||Risk Monitoring
188||12|0|Risk Reviews: Risks will be reviewed during weekly progress meetings to assess their status and implement mitigation strategies.
189||12|0|Issue Tracking: All risks and issues will be logged in the project management tool Jira and assigned to team members for resolution.
190|heading 4|||Reporting Frequency
191||8|0|Daily Standups: Quick updates on individual progress, blockers, and plans for the day.
192||8|0|Weekly Reports: Summary of completed tasks, risks, and upcoming priorities.
193||8|0|Final Project Report: A comprehensive report at the end of the project, detailing deliverables, lessons learned, and overall performance.
203|heading 2|||9. Conclusion
204||||This initial plan outlines the project structure, scheduling, and risks, providing a comprehensive roadmap for the successful development and deployment of the Schedify: Constraint-Based Automated Exam Scheduling System. By clearly defining the scope, deliverables, and resource requirements, the plan ensures alignment with stakeholder expectations and establishes a solid foundation for execution.
205||||The project team has identified potential risks and mitigation strategies, ensuring that challenges such as requirement changes, developer turnover, and performance bottlenecks are proactively addressed. The use of Agile methodology, regular stakeholder engagement, and robust monitoring tools like Jira and GitHub will enable the team to adapt to changes and maintain progress toward milestones.
206||||With a well-defined work breakdown structure, realistic timelines, and a focus on quality assurance, the project is positioned to deliver a scalable, efficient, and user-friendly solution that meets the needs of university departments. Regular updates, sprint reviews, and stakeholder feedback will ensure continuous improvement and alignment with project goals.
207||||The journey to success begins with a single step, and this plan is ours. With clarity, collaboration, and commitment, Schedify will achieve its goal of transforming exam scheduling."